
\vspace*{0.2cm}

{\noindent\textbf{\nameref{\en{EN}kapEinleitung}~~\dotfill~~\pageref{\en{EN}kapEinleitung}}\par\vspace*{8pt}}

\kommentInh{ellipfkt}{\en{We develop the terms and propositions about the Weierstraß elliptic functions that we need for our proof of the Chudnovsky formula.}%
\de{Alle für die Herleitung der Chudnovsky-Formel benötigten Begriffe und Sätze über die Weierstraß'schen elliptischen Funktionen werden entwickelt.}}

\kommentInh{kapquasiint}{\en{We define the quasiperiods of a lattice with the Weierstraß $\zeta$-function. Then we give an alternative representation of the periods and quasiperiods with means of elliptic integrals.}%
\de{Die \glqq Quasiperioden\grqq~eines Gitters werden mit Hilfe der Weierstraß'schen $\zeta$-Funk"-tion definiert. Es folgt eine alternative Darstellung der Perioden und Quasiperioden mit Hilfe elliptischer Integrale.}}

\kommentInh{kapGitter}{\en{In this chapter, we will see that two lattices that are rotated and/or scaled versions of each other can be called "equivalent" and that equivalent lattices have the same value of Klein's absolute invariant~$J$.}%
\de{In diesem Kapitel werden wir sehen, dass zwei Gitter, die durch eine Drehstreckung auseinander hervorgehen, \glqq äquivalent\grqq~genannt werden können und dass äquivalente Gitter die gleiche absolute Invariante $J$ haben.}}

\kommentInh{kapFourier}{\en{We calculate the Fourier representations of the normalized Eisenstein series.}%
\de{Die Fourierentwicklungen der normierten Eisensteinreihen werden bewiesen.}}

\kommentInh{kapanhang}{\en{These estimates prove that Kummer's solution in chapter~\ref{\en{EN}kapkummer} converges, and they are needed to calculate the coefficients in chapter~\ref{\en{EN}kapFormel}.}%
\de{Die hier bewiesenen Abschätzungen garantieren, dass Kummers Lösung in Kap.~\ref{\en{EN}kapkummer} konvergiert und ermöglichen die Berechnungen der Koeffizienten in Kap.~\ref{\en{EN}kapFormel}.}}

\kommentInh{kapClausen}{\en{We prove Clausen's formula and the hypergeometric differential equations needed for the proof. This chapter is self-contained.}%
\de{Wir beweisen die Clausen-Formel und die für den Beweis benötigten hypergeometrischen Differentialgleichungen. Das Kapitel ist unabhängig von den vorigen.}}

\kommentInh{kappicardfuchs}{\en{This proof of the Picard Fuchs differential equation can be read straight after chapter~\ref{\en{EN}kapGitter}.}%
\de{Dieser Beweis der Picard-Fuchs-Differentialgleichung kann direkt im Anschluss an Kap.~\ref{\en{EN}kapGitter} gelesen werden.}}

\kommentInh{kapkummer}{\en{We use one of Kummer's solutions of the Picard Fuchs differential equation to prove a connection between the periods of a lattice and a hypergeometric function.}%
\de{Mit Hilfe einer Kummer'schen Lösung der Picard-Fuchs-Differentialgleichung wird ein Zusammenhang zwischen den Perioden eines Gitters und einer hypergeometrischen Funktion hergestellt.}}

\kommentInh{kaphaupt}{\en{We prove the Main Theorem~\ref{\en{EN}hauptformel} using Kummer's solution, Clausen's formula and the Fourier representations.}%
\de{Das Haupttheorem~\ref{\en{EN}hauptformel} wird bewiesen, ausgehend von Kum"-mers Lö"-sung und mit Hilfe der Clausen-Formel und der Fourierdarstellungen.}}

\kommentInh{kapFormel}{\en{We explicitly calculate the exact values of $s_2(\tau_{N})$ and $J(\tau_{N})$ using the estimates from chapter~\ref{\en{EN}kapanhang}. Thus we obtain the Chudnovsky formula and ten further formulae to calculate $\pi$.}%
\de{Die exakten Werte von $s_2(\tau_{N})$ und $J(\tau_{N})$ werden mit Hilfe der Abschätzungen aus Kap.~\ref{\en{EN}kapanhang} explizit berechnet. Somit erhalten wir die Chudnovsky-Formel und zehn weitere Formeln zur Berechnung von $\pi$.}}

\kommentInh{kapdivi}{\en{We prove that $m\cdot\wp(u;L)$ is an algebraic integer of~~$\mathbb Z\mathopen{}\left[\frac 1 4 g_2(L); \frac 1 4 g_3(L)\right]\mathclose{}$ for all positive integers $m$ and for all $u\in\mathbb C-L$ with $m\cdot u\in L$.}%
\de{Wir beweisen, dass $m\cdot\wp(u;L)$ für alle natürlichen Zahlen $m\neq 0$ und für alle $u\in\mathbb C-L$ mit $m\cdot u\in L$ ganzalgebraisch in $\mathbb Z\mathopen{}\left[\frac 1 4 g_2(L); \frac 1 4 g_3(L)\right]\mathclose{}$ ist.}}

\kommentInh{kapmult}{\en{We use Appendix~\ref{\en{EN}kapdivi} to prove that $\sqrt{D}\cdot\frac{E_2^*(\tau)}{\eta^4(\tau)}\cdot(AC)^2$ is an algebraic integer if $\tau$ satisfies $C\tau^2+B\tau+A=0$ with discriminant $D$.}%
\de{Wir beweisen mit Hilfe von Anhang~\ref{\en{EN}kapdivi}, dass $\sqrt{D}\cdot\frac{E_2^*(\tau)}{\eta^4(\tau)}\cdot(AC)^2$ ganzalgebraisch ist, falls $\tau$ eine Lösung von $C\tau^2+B\tau+A=0$ mit Diskriminante $D$ ist.}}

{\noindent\textbf{\hyperref[\en{EN}literat]{\en{References}\de{Literatur}}~~\dotfill~~\pageref{\en{EN}literat}}}


